\begin{table}[t]
  \centering
  \caption{\textbf{座標フレーム。}
  すべての剛体変換は正規直交・行列式 \(+1\) の同次行列として保存する。
  一様尺度はメートルアダプタが持ち、
  剛体変換の内部へは隠さない。
  }
  \label{tab:frames}
  \small
  \begin{tabularx}{\linewidth}{@{}p{0.20\linewidth}p{0.24\linewidth}X@{}}
    \toprule
    フレーム & 単位と軸 & 契約 \\
    \midrule
    正規化シーン \(\Scene_n\) & 正規化、右手系 &
      再構成が公開する座標。ガウス平均と復元カメラはここに住む。 \\
    メートルシーン \(\Scene_m\) & メートル、右手系 &
      \(\Scene_n\) の一様尺度変換。受理された全コートの共通の土台。 \\
    コート \(\Court_i\) & メートル。\(+X\) 右サイドライン、
      \(+Y\) 遠ベースライン、\(+Z\) 上 &
      規制寸法と物理点の番号付け。 \\
    カメラ \(\Cam\) & OpenCV。\(+x\) 右、\(+y\) 下、\(+z\) 前 &
      \(\T{\Scene_m}{\Cam}\) はカメラからシーン。
      投影はその逆変換を用いる。 \\
    正準コート \(\Canon\) & メートル、対象コート相対 &
      カメラ端の選択 \(S\in\{I,R_z(\pi)\}\) を適用した後のコートフレーム。
      姿勢教師の基準。 \\
    モデル画像 & 補正後画素 &
      キーポイントのロジット、主点、焦点距離、再投影残差を同じフレームで扱う。 \\
    \bottomrule
  \end{tabularx}
\end{table}
